\chapter{Literature Review}
\label{ch:literature_review}

\section{Overview of the Chapter and Research Context}
\label{sec:lit_chapter_overview}
When I started doing research on input interfaces, I came to understand that most human-computer interactions still make use of key switches. The use of physical keyboards provides reliable tactile feedback for typing prose, but modern software and creative jobs need macro keypads and shortcut consoles \cite{ReviewVirtualKeyboard2020, Lee2022Virtual}. These include video editing consoles and developer launching pads. But commercially available hardware macro pads cost from \$100 to \$300, are not customizable in any way physically, and it is hard to maintain hygiene in them. This led researchers to look into the possibility of using virtual keyboards on flat surfaces by making use of computer vision \cite{Zhang2001Visual, Srivastava2012RealTime, Khare2019QWERTY, Maman2023TypeNet}. Although most early attempts have been made to create alternatives for continuous typing on paper, it was physically difficult to type on paper continuously without any tactile switch travel. For customizable macro controllers, paper works much better because users only glance at and tap individual shortcut buttons.

There were four huge challenges to getting a monocular paper virtual keyboard to work well:
\begin{itemize}
    \item \textbf{Surface Touch Detection Without Depth Sensors:} With standard 2D webcams it is impossible to measure depth in the Z-axis, making it difficult to differentiate between a real paper touch and a hovering finger \cite{Thomas2013Camera, Posner2012Single, Lee2019Virtual, Toshpulatov2024RealTime}.
    \item \textbf{Hand Scale and Perspective Variations:} Hand landmark coordinates vary according to different hand sizes and distances from the lens \cite{Zhang2020MediaPipe, Doan2022Efficient}. Normalizing coordinates is needed without adding lag from heavy smoothing filters \cite{Kumar2026Hybrid}.
    \item \textbf{Planar Orientation and Camera Tilt:} The perspective effect on the desk surface that comes from tilting the camera means we need to compute a planar homography matrix ($H \in \mathbb{R}^{3 \times 3}$) so as to keep key bounds aligned continuously \cite{Wang2016AprilTag2, Kallwies2020Determining, Pirchheim2011Homography}.
    \item \textbf{Decoupled Layout Design and Action Multiplexing:} Older systems hardcoded QWERTY layouts rather than allowing users to pair various software action profiles with one printed sheet \cite{ReviewVirtualKeyboard2020}.
\end{itemize}

In order to lay a sound foundation, this chapter analyzes the current literature across four main categories. Firstly, I define the literature framework and searching methodology. Secondly, I discuss the evolution of input devices and the biomechanical aspects of finger strikes. Thirdly, I conduct a comparative analysis of existing virtual keyboards from eight different sensing mechanisms and analyze their drawbacks. Lastly, I examine the mathematical and algorithmic aspects of MediaPipe hand tracking, scale normalization, AprilTag homography, and PyTorch LSTM touch detection, leading to the research gaps that motivate this thesis.

\section{Conceptual Map of the Literature}
\label{sec:lit_conceptual_map}
In Figure~\ref{fig:lit_conceptual_map_diagram}, I plotted the literature against four primary pillars: HCI input approaches, sensing types, technical algorithms, and synthesis of research gaps.

\begin{figure}[htbp]
\centering

\includegraphics[width=0.85\textwidth]{figures/chapter03_fig01.png}

\caption{Conceptual taxonomy and structural organization of the literature review.}
\label{fig:lit_conceptual_map_diagram}
\end{figure}

The review is structured across four main themes:
\begin{itemize}
    \item \textbf{HCI and Input Paradigms (Section~\ref{sec:lit_domain_overview}):} Evolution of input hardware, biomechanics of finger-surface impact, and user performance metrics.
    \item \textbf{Sensing Modalities (Section~\ref{sec:lit_existing_systems}):} Critical evaluation of eight sensing technologies, synthesized in Table~\ref{tab:lit_comparative_matrix}.
    \item \textbf{Technological Analysis (Section~\ref{sec:lit_technological_analysis}):} Mathematical formulation and algorithmic analysis of hand tracking, scale normalization, fiducial homography, and deep learning touch detection.
    \item \textbf{Research Synthesis (Section~\ref{sec:lit_reflection}):} Identification of the six fundamental literature gaps and the formal research problem statement.
\end{itemize}

\subsection{Literature Search Strategy and PRISMA 2020 Systematic Protocol}
\label{subsec:lit_search_strategy}
For the systematic review process, I applied the PRISMA 2020 protocol through IEEE Xplore, ACM Digital Library, Scopus, Elsevier ScienceDirect, and Google Scholar to find peer-reviewed studies from 2000 to 2025 \cite{ReviewVirtualKeyboard2020}:
\begin{enumerate}
    \item \textbf{Search Syntax:} I searched with queries containing: \texttt{("virtual keyboard" OR "paper keyboard") AND ("monocular RGB" OR "MediaPipe") AND ("touch detection" OR "homography")}.
    \item \textbf{Inclusion Criteria:} I added peer-reviewed papers that showed functioning virtual keyboards, planar contact algorithms, fiducial tracking with tilt, or real-time hand joint estimation.
    \item \textbf{Exclusion Criteria:} I removed papers that required wearable sensor gloves, intrusive markers, multi-camera motion capture rigs, or were only conceptual white papers without implementation.
    \item \textbf{Synthesis Approach:} The shortlisted papers were then compared on the basis of sensing approaches, multi-finger support, compute latency, and environmental robustness to identify the problem areas this project addresses.
\end{enumerate}

\section{Domain Overview}
\label{sec:lit_domain_overview}

\subsection{Evolution of Text Entry and Spatial HCI Hardware}
\label{subsec:lit_history_evolution}
Considering the evolution of text input, early mechanical keyboards used physical switches that provided key travel with strong tactile feedback \cite{ReviewVirtualKeyboard2020}. Next came membrane keyboards, where production became cheaper, and capacitive touchscreens that offered soft visual key layouts without physical tactile edges \cite{ReviewVirtualKeyboard2020}. Optical laser keyboards projected the outline of keys using infrared light \cite{Cheng2015Fingertip, Kudale2016RealTime}; however, these could not work reliably in bright environments as ambient light would wash out the pattern. Nowadays, entering complex command sequences can be done by dedicated macro pads (such as the Stream Deck), which enable users to trigger multi-key shortcuts with one keypress. However, such devices are expensive (usually priced from \$100 up to \$300) while being non-modifiable physically. Monocular computer vision allows us to track 21 hand joints using only a standard webcam \cite{Zhang2020MediaPipe, GilMartin2023Hand}, which means we can use a printed piece of paper as an inexpensive, customizable macro pad without active electronics.

\subsection{Kinematic Typing Dynamics and Biomechanical Surface Impact}
\label{subsec:lit_typing_biomechanics}
Without a depth sensor, I needed to understand the physical characteristics of a finger as it makes contact with a flat desk in order to identify touch events \cite{Simao2016Unsupervised, MacRitchie2015Integrating, Pilkov2024Markerless}. A physical keystroke consists of four kinematic phases:
\begin{enumerate}
    \item \textbf{Approaching Phase ($t_1$ to $t_2$):} The finger moves downward at increasing vertical speed towards the target button ($V_y$).
    \item \textbf{Impact Phase ($t_2$ to $t_3$):} When encountering the rigid paper surface, downward velocity is zeroed in 20 to 40~ms ($V_y \to 0$), producing a steep deceleration spike ($A_y = \frac{dV_y}{dt} \ll 0$).
    \item \textbf{Resting Phase ($t_3$ to $t_4$):} The finger pauses briefly on the paper while the touch event registers.
    \item \textbf{Recoil Phase ($t_4$ to $t_5$):} The finger rebounds upward away from the desk surface ($V_y > 0$).
\end{enumerate}
The rigid desk acts as a hard physical boundary, providing immediate deceleration of motion. Conversely, fingers tend to overshoot in mid-air typing systems, as there is no stopping surface \cite{Boletsis2019TextInput, Maman2023TypeNet}.

\subsection{Ergonomics and Human Input Performance Metrics}
\label{subsec:lit_fitts_law}
Input performance by humans has a direct correlation with physical feedback and movement dynamics:
\begin{itemize}
    \item \textbf{Tactile Key Travel versus Discrete Visual Targeting:} With a mechanical keyboard, a depth of keystroke between $2\text{ to }4\text{ mm}$ allows fast blind touch-typing without needing to see the keys. On flat paper, typing continuous prose creates strain on fingers since there is no mechanical key travel \cite{ReviewVirtualKeyboard2020}. For macro keypads, people do not type continuous prose; rather, they look for a particular button and press it once. The presence of custom button sizes on paper layouts directly improves ergonomics.
    \item \textbf{Discrete Macro Action Rate versus Continuous Typing Speed:} Traditional typing tests measure input speed in Words Per Minute based on continuous character sequences. In contrast, for macro keypads, touch accuracy and reduced latency are the critical performance metrics rather than continuous typing speed.
    \item \textbf{Keystrokes Per Character (KSPC):} Efficiency is measured by Keystrokes Per Character. Physical keyboards generally register about 1.05 KSPC, while contactless mid-air keyboards have more than 1.30 KSPC owing to accidental touches \cite{Boletsis2019TextInput}.
    \item \textbf{Fitts' Law:} Fitts' Law models how long it takes to reach a target button of width $W$ at distance $D$ through the Index of Difficulty ($ID$) in bits:
    \begin{equation}
    MT = a + b \log_2 \left( 1 + \frac{D}{W} \right) = a + b \cdot ID
    \label{eq:fitts_law}
    \end{equation}
    The customizable paper macro keyboard enables one to customize the button width $W$ in the software designer, directly reducing the Index of Difficulty for rapid shortcut selection.
\end{itemize}

\subsection{Multi-Finger Tendon Coupling}
\label{subsec:lit_biomechanical_coupling}
Anatomical tendon coupling in the human hand is another biological factor \cite{Simao2016Unsupervised, Pilkov2024Markerless}. Shared muscle pathways are found in the flexor and extensor tendons of the middle, ring, and pinky fingers. For example, when someone presses down with one finger, the fingers next to it involuntarily twitch downward. These minor twitches are confused for intentional key presses by basic single-finger tracking algorithms \cite{Thomas2013Camera, Srivastava2012RealTime}. Deep sequence models address this by tracking combinations of movement from all 21 hand joints over time to differentiate intentional taps from passive tendon twitches.

\section{Existing Systems and Sensing Modalities}
\label{sec:lit_existing_systems}
Reviewing previous systems of virtual keyboards, one may note that eight different sensing modalities have been explored.

\subsection{Projection-Based and Infrared Sensor Systems}
\label{subsec:lit_proj_systems}
Projection keyboards use a red laser grid placed on the table surface and sense finger presence as an invisible infrared beam is broken \cite{Cheng2015Fingertip, Kudale2016RealTime}. Kudale and Wanjale \cite{Kudale2016RealTime} implemented an infrared camera tracking reflective markers using four-point homography. However, the ambient light of the room destroys the pattern through wash-out and makes changing the key arrangement impossible.

\subsection{Shadow-Based Monocular RGB Keyboards}
\label{subsec:lit_shadow_systems}
Monocular RGB shadow sensors calculate the distance between the fingertip and its cast shadow, which reaches zero upon physical contact \cite{Thomas2013Camera, Posner2012Single, Yue2014Blind}. Posner et al. \cite{Posner2012Single} applied adaptive shadow thresholding for paper contact detection. However, shadow analysis fails under diffuse lighting, multiple light sources, or dark desk surfaces \cite{Yue2014Blind}. Crucially, shadow methods completely fail when several fingers move simultaneously because their shadows overlap and merge.

\subsection{Hardware-Assisted Depth Sensors (RGB-D and ToF)}
\label{subsec:lit_depth_systems}
Depth sensors like Microsoft Kinect, Intel RealSense, and Time-of-Flight cameras create metric 3D point clouds \cite{Lee2019Virtual, Toshpulatov2024RealTime}. Lee and Kwon \cite{Lee2019Virtual} calibrated a flat surface to detect touches when fingertip point clouds crossed a distance threshold. Toshpulatov et al. \cite{Toshpulatov2024RealTime} combined 3D depth information and optical flow to filter out hovering artifacts. However, depth cameras are costly, power-hungry, prone to noise close to flat surfaces, and require dedicated GPUs to process in real-time.

\subsection{3D Mid-Air Typing and AR/VR Virtual Keyboards}
\label{subsec:lit_air_systems}
Spatial computing headsets have spurred interest in mid-air keyboards \cite{Boletsis2019TextInput, Enkhbat2020HandKey, Lee2022Virtual, Yoo2026WordLevel}. Enkhbat et al. \cite{Enkhbat2020HandKey} introduced HandKey, detecting typing within 3D bounding cubes. Lee et al. \cite{Lee2022Virtual} tracked two-handed mid-air typing for head-mounted displays. However, mid-air typing lacks a physical desk surface where fingers stop; therefore, users experience high levels of arm fatigue and overshoot target keys \cite{Boletsis2019TextInput, Maman2023TypeNet}.

\subsection{Paper-Anchored Monocular RGB Keyboards}
\label{subsec:lit_paper_systems}
Monocular paper keyboards employ standard webcams to interact virtually with printed paper sheets \cite{Zhang2001Visual, Srivastava2012RealTime, Khare2019QWERTY, Maman2023TypeNet}. Zhang et al. \cite{Zhang2001Visual} tracked page corners to compute planar homography. Vertical velocity thresholds of pixels were employed by Srivastava and Tripathi \cite{Srivastava2012RealTime}, while a printed QWERTY sheet was tracked by Khare \cite{Khare2019QWERTY} employing skin-color segmentation. TypeNet \cite{Maman2023TypeNet} enabled typing on uncalibrated surfaces with deep keypoint tracking. However, all prior monocular paper keyboards were limited in that they detected the movement of only one finger at a time.

\subsection{Specialized and Assistive Modalities}
\label{subsec:lit_specialized_assistive}
Additional forms of sensing involve VLC photodiodes \cite{Webber2021Finger}, fingernail color-based blood volume imaging \cite{Sun2009Estimation}, optical motion capture \cite{MacRitchie2015Integrating, Pilkov2024Markerless}, and gaze tracking \cite{Chen2024Gaze}. Gaze tracking assists people with motor disabilities, although the Midas Touch effect causes unintended key clicks while stopping to read \cite{Chen2024Gaze}.

\subsection{Fiducial Tracking in Planar Robotics}
\label{subsec:lit_robotics_systems}
With regard to tracking technology, planar robotics research proves that square fiducial markers like AprilTag and ArUco offer sub-pixel tracking accuracy even when the camera is tilted severely \cite{Ivacko2024Edge, Herdiansyah2024Implementation, Schenck2024Utilizing, GarridoJurado2026ArUco}. Ivacko et al. \cite{Ivacko2024Edge} utilized ArUco markers for mapping robot coordinate frames, while Garrido-Jurado et al. \cite{GarridoJurado2026ArUco} realized high frame rate contour scanning.

\subsection{Comparative Literature Analysis Matrix}
\label{subsec:lit_comparative_matrix}
Some examples of previous work in terms of sensing modality, hardware requirements, tracking techniques, scale normalization, temporal filtering, touch classification, and layout flexibility are presented in Table~\ref{tab:lit_comparative_matrix}. Importantly, across all monocular RGB-based approaches operating on physical surfaces, none support multi-finger interaction; multi-finger input was previously available only in expensive depth camera setups or virtual reality environments.

\begin{table}[htbp]
\centering
\caption{Comparative analysis of existing virtual keyboard frameworks and sensing methodologies.}
\label{tab:lit_comparative_matrix}
\begin{tabular}{l l l l l l l l l}
\toprule
\textbf{Study and Year} & \textbf{Sensing Modality} & \textbf{Hardware Req.} & \textbf{Surface Tracking} & \textbf{Scale Normalization} & \textbf{Temporal Filter} & \textbf{Touch Classifier} & \textbf{Latency / FPS} & \textbf{Layout Customization} \\
\midrule
Zhang et al. (2001) \cite{Zhang2001Visual} & Monocular RGB & Standard Webcam & Quad Paper Bounds & None (Pixel Coords) & Moving Avg. & Spatial Proximity & High ($\sim 15$ FPS) & Rigid Paper Quad \\
Sun et al. (2009) \cite{Sun2009Estimation} & Monocular RGB & High-Res Camera & None (Fixed Finger) & None & None & LDA Nail Coloration & Offline ($< 5$ FPS) & None \\
Pirchheim and Reitmayr (2011) \cite{Pirchheim2011Homography} & Monocular RGB & Mobile Camera & Keyframe Homography & Feature Scale & Particle Filter & Spatial Tracking & Real-Time ($30$ FPS) & None (SLAM Map) \\
Srivastava and Tripathi (2012) \cite{Srivastava2012RealTime} & Monocular RGB & Standard Webcam & Manual Quad Contour & None (Raw Pixels) & Static Smoothing & Speed Deceleration & Real-Time ($25$ FPS) & Manual Paper Drawing \\
Posner et al. (2012) \cite{Posner2012Single} & Monocular RGB & Standard Webcam & Fixed Desk Layout & None & Dynamic Threshold & Fingertip-Shadow Dist. & Real-Time ($30$ FPS) & Fixed Paper Template \\
Thomas (2013) \cite{Thomas2013Camera} & Monocular RGB & Standard Webcam & Fixed Area & None & None & Shadow Gradient Vector & Real-Time ($30$ FPS) & Fixed Layout \\
Yue et al. (2014) \cite{Yue2014Blind} & Monocular RGB & Smartphone RGB & Screen Border H & None & Optical Flow & Shadow + k-Means & Offline Analysis & Fixed Smartphone Screen \\
Cheng et al. (2015) \cite{Cheng2015Fingertip} & Proj. + Camera & Projector + Camera & Photometric Calib. & None & Background Sub. & 3D Shadow Distance & Real-Time ($20$ FPS) & Fixed Projected GUI \\
MacRitchie and McPherson (2015) \cite{MacRitchie2015Integrating} & Optical MoCap & MoCap + Touch & Surface Calibration & 3D Skeleton mm & Moving Average & Velocity Inversion & Real-Time ($100$ FPS) & Fixed Piano Keys \\
Kudale and Wanjale (2016) \cite{Kudale2016RealTime} & IR Webcam & IR Filtered Camera & 4-Point Homography & None & Thresholding & IR Pen Tap Contact & Real-Time ($30$ FPS) & Projected Display \\
Simao et al. (2017) \cite{Simao2016Unsupervised} & Motion Capture & Leap Motion Sensor & None & Bone Lengths & Sliding Window & Velocity Inversion & Real-Time ($60$ FPS) & None \\
Ji et al. (2018) \cite{Ji2018Local} & Monocular RGB & Standard Webcam & Skin Color Mask & None & Moving Avg. & k-Means Clustering & Real-Time ($25$ FPS) & Dynamic Key Clusters \\
Nunez et al. (2018) \cite{Nunez2018Convolutional} & 3D Skeleton & Depth Sensor & None & Normalized Skeleton & Sequence Window & CNN + LSTM Network & Real-Time ($30$ FPS) & None \\
Lai and Yanushkevich (2018) \cite{Lai2018CNN} & Multimodal RGB-D & Depth Sensor & None & Skeleton Joint Map & Temporal Buffer & CNN + RNN Fusion & Real-Time ($30$ FPS) & Gesture Set \\
Khare (2019) \cite{Khare2019QWERTY} & Monocular RGB & Standard Webcam & Static Bounds & None & Color Threshold & Pause Over Bounds & Real-Time ($20$ FPS) & Fixed QWERTY Paper \\
Lee and Kwon (2019) \cite{Lee2019Virtual} & RGB-D Depth & ToF Depth Camera & Depth Plane Calib. & 3D Metric mm & Weight Filtering & Z-Depth Threshold & Real-Time ($30$ FPS) & Fixed Virtual Touch \\
Boletsis and Kongsvik (2019) \cite{Boletsis2019TextInput} & VR Controller & VR Headset & Virtual 3D Grid & VR World Scale & None & Striking Velocity & Real-Time ($90$ FPS) & Fixed VR Layout \\
Enkhbat et al. (2020) \cite{Enkhbat2020HandKey} & Monocular RGB & Standard Webcam & Air Typing Zone & Hand Size Scaling & None & CNN Spatial Classifier & Real-Time ($30$ FPS) & Fixed Air QWERTY \\
Cui et al. (2020) \cite{Cui2020Deep} & Vision Review & Various Sensors & Various & Various & Various & Deep Learning Survey & N/A & N/A \\
Zhang et al. (2020) \cite{Zhang2020MediaPipe} & Monocular RGB & Commodity Camera & Palm Detector & Normalized 21 Keypoints & Single-frame Regress & Skeletal Keypoint Reg. & High Speed ($> 60$ FPS) & General Hand Tracking \\
Ikram and Liu (2020) \cite{Ikram2020Skeleton} & 3D Skeleton & Leap Motion Sensor & None & Hand Scale & Sliding Window & 1D CNN + LSTM & Real-Time ($50$ FPS) & Mid-air Gestures \\
Wang et al. (2020) \cite{Wang2020Relative} & Monocular RGB & Stereo Pair & Superpixel Homography & Superpixel Scale & Bundle Adjust & Multi-Model RANSAC & Real-Time ($25$ FPS) & Planar Maps \\
Li et al. (2021) \cite{Li2021CNN} & Monocular RGB & Embedded Camera & Grid Mask & None & CNN Feature Map & Lightweight CNN & Edge Real-Time ($30$ FPS) & Fixed Key Grid \\
Webber et al. (2021) \cite{Webber2021Finger} & Visible Light & VLC Photodiodes & Optical Intensity Mask & None & Intensity Features & SVM / Random Forest & Real-Time ($30$ FPS) & Light Grid Keys \\
Gammulle et al. (2021) \cite{Gammulle2021TMMF} & Multimodal RGB & Single Camera & None & Sequence Normalization & Multi-Modal Fusion & Single-Stage TMMF Net & Real-Time ($30$ FPS) & Continuous Gestures \\
Zhang et al. (2022) \cite{Zhang2022DynaKey} & Head-Mounted RGB & Smart Glasses & Gyro + Homography & None & Motion Compensation & Downward Trajectory & Real-Time ($30$ FPS) & Drawn Keyboard Bounds \\
Lee et al. (2022) \cite{Lee2022Virtual} & RGB-D / HMD & AR/VR Headset & HMD Frame & 3D Joint Scaling & Temporal Buffer & 7-Class Deep Network & Real-Time ($60$ FPS) & Ambidextrous Layouts \\
Doan (2022) \cite{Doan2022Efficient} & Monocular RGB & Commodity RGB & None & Scale Normalized & Coordinate Smooth & Multi-layer LSTM & Real-Time ($30$ FPS) & Body Activity \\
Yang et al. (2022) \cite{Yang2022Towards} & Smartphone RGB & Commodity Phone & Uncalibrated Video & Self-Supervised Scale & Motion Trajectory & Deep Inference Model & Offline Analysis & Uncalibrated Keyboards \\
Maman and Bermano (2023) \cite{Maman2023TypeNet} & Monocular RGB & Front Camera & Uncalibrated Plane & Hand Skeleton Scale & Self-Refinement & Deep Net + LM Decoding & Real-Time ($30$ FPS) & Fixed QWERTY Flat \\
Gil-Martin et al. (2023) \cite{GilMartin2023Hand} & Monocular RGB & Commodity Camera & None & MediaPipe Scale & Sequence Buffer & Deep Recurrent Net & Real-Time ($60$ FPS) & General Gestures \\
Sanchez-Brizuela et al. (2023) \cite{SanchezBrizuela2023Lightweight} & Monocular RGB & Commodity Camera & None & Landmark Distance & Morphological Mask & Landmark Rule Classifier & Ultra High ($90$ FPS) & Skin Segmentation \\
Ivacko et al. (2024) \cite{Ivacko2024Edge} & Monocular RGB & Wrist Camera & ArUco Calibration & Metric Scaling & Low-Pass Alpha & Rule-Based Kinematics & Edge Real-Time ($30$ FPS) & Robot Control Pad \\
Herdiansyah et al. (2024) \cite{Herdiansyah2024Implementation} & Monocular RGB & Single RGB Camera & ArUco PnP & Metric Tag Size & Zhang Calibration & Geometric Pose Engine & Real-Time ($30$ FPS) & Navigation Targets \\
Pilkov (2024) \cite{Pilkov2024Markerless} & Multi-Camera RGB & 3 RGB Cameras & None & 3D Skeleton mm & Temporal Alignment & Deep Keypoint Trajectory & Real-Time ($30$ FPS) & Piano Keyboard \\
Toshpulatov et al. (2024) \cite{Toshpulatov2024RealTime} & Multimodal RGB-D & Depth + RGB & Surface Calib. & 3D Spatial mm & Fusion Filter & Temporal Depth-Motion Net & Real-Time ($30$ FPS) & VR/AR GUI Objects \\
Chen (2024) \cite{Chen2024Gaze} & Monocular RGB & Commodity Webcam & Screen Grid & Pupil Vector Scale & Dwell Time Filter & Eye-Gaze Deep Model & Real-Time ($30$ FPS) & Gaze Keypad \\
Schenck et al. (2024) \cite{Schenck2024Utilizing} & Monocular RGB & Single RGB Camera & ArUco Inpainting & Metric Scale & Inpainting Mask & Deep Keypoint Detector & Real-Time ($30$ FPS) & Robotic Links \\
Andriyanov et al. (2025) \cite{Andriyanov2025Improving} & Monocular RGB & Commodity RGB & Bounding Box & Normalized Bounds & Temporal Window & MediaPipe + YOLO-Pose & Real-Time ($45$ FPS) & Gesture Set \\
Yoo and Lee (2026) \cite{Yoo2026WordLevel} & Monocular RGB & Single RGB Camera & None & Trajectory Scaling & Displacement Accum. & Word Motion Network & Real-Time ($30$ FPS) & Contactless QWERTY \\
Kumar (2026) \cite{Kumar2026Hybrid} & Monocular RGB & Commodity Camera & Canvas Plane & Landmark Ratio & Adaptive Low-Pass & Motion Threshold & Real-Time ($60$ FPS) & Air Canvas Drawing \\
Garrido-Jurado et al. (2026) \cite{GarridoJurado2026ArUco} & Monocular RGB & Commodity Camera & ArUco Nano Tags & Metric Tag Scale & Sub-pixel Sampling & Visited-Aware Contours & Edge High ($> 100$ FPS) & Robotics Targets \\
\bottomrule
\end{tabular}%
\end{table}

\section{Technological Analysis}
\label{sec:lit_technological_analysis}
This section analyzes the four foundational technical modules: skeletal hand tracking, scale normalization, planar homography, and deep learning touch classification.

\subsection{Hand Pose Estimation and Skeletal Joint Tracking}
\label{subsec:lit_hand_pose}
Markerless hand tracking is a must for natural interactions without sensors on the hand \cite{Cui2020Deep, Zhang2020MediaPipe}. Initially, color thresholding was attempted in HSV space \cite{Ji2018Local, Khare2019QWERTY}, but this failed as light conditions changed or due to differences in skin tones. Google MediaPipe Hands \cite{Zhang2020MediaPipe} resolves this problem with the help of a two-stage deep learning process optimized to work well on regular CPUs:
\begin{enumerate}
    \item \textbf{Single-Shot Palm Detector:} A one-shot palm detector finds a bounding box for the palm region across the entire image using focal loss $\mathcal{L}_{\text{focal}} = -\alpha_t (1 - p_t)^\gamma \log(p_t)$. By focusing on the rigid palm region rather than articulating finger movement, tracking loss is minimized.
    \item \textbf{Hand Landmark Model:} The cropped palm region feeds into a hand landmark model that predicts 3D coordinates $(x_i, y_i, z_i)$ for 21 skeletal joints (Figure~\ref{fig:mediapipe_hand_topology}).
\end{enumerate}

\begin{figure}[htbp]
\centering

\includegraphics[width=0.85\textwidth]{figures/chapter03_fig02.png}

\caption{MediaPipe 21 hand landmark anatomical topology.}
\label{fig:mediapipe_hand_topology}
\end{figure}

MediaPipe provides five fingertip coordinates (Landmarks 4, 8, 12, 16, 20) and knuckle base joints (MCP, PIP, DIP) with execution speeds greater than 60~FPS on standard CPUs \cite{GilMartin2023Hand, SanchezBrizuela2023Lightweight}.

\subsection{Kinematic Scale Normalization and Raw Feature Propagation}
\label{subsec:lit_scale_norm}
Landmark pixel positions vary with the distance between the user's hand and the webcam, along with the actual size of the hand \cite{Doan2022Efficient, Kumar2026Hybrid}. With respect to the standard pinhole camera model of projection:
\begin{equation}
x_{\text{pixel}} = f \frac{X}{Z}, \quad y_{\text{pixel}} = f \frac{Y}{Z}
\label{eq:pinhole_projection}
\end{equation}

Scaling camera distance by factor $k$ ($Z' = k Z$) implies inverse scaling of pixel positions:
\begin{equation}
x'_{\text{pixel}} = \frac{x_{\text{pixel}}}{k}, \quad y'_{\text{pixel}} = \frac{y_{\text{pixel}}}{k}, \quad L'_{\text{hand}} = \frac{L_{\text{hand}}}{k}
\label{eq:pixel_scaling}
\end{equation}
where $L_{\text{hand}}$ is the reference hand length defined as the Euclidean distance between the wrist $\mathbf{P}_0 = (x_0, y_0)$ and the middle MCP knuckle $\mathbf{P}_9 = (x_9, y_9)$:
\begin{equation}
L_{\text{hand}} = \sqrt{(x_9 - x_0)^2 + (y_9 - y_0)^2}
\label{eq:hand_length_scale_def}
\end{equation}

Scaling landmark displacements with respect to the wrist reference $\mathbf{P}_0$ by hand length $L_{\text{hand}}$ eliminates the influence of scale variations:
\begin{equation}
\mathbf{P}_i^{\text{norm}\prime} = \frac{\mathbf{P}_i' - \mathbf{P}_0'}{L'_{\text{hand}}} = \frac{\frac{\mathbf{P}_i - \mathbf{P}_0}{k}}{\frac{L_{\text{hand}}}{k}} = \frac{\mathbf{P}_i - \mathbf{P}_0}{L_{\text{hand}}} = \mathbf{P}_i^{\text{norm}}
\label{eq:scale_invariance_proof}
\end{equation}

First-order numerical velocities are calculated between consecutive frames:
\begin{equation}
\mathbf{V}_i(t) = \frac{\mathbf{P}_i^{\text{norm}}(t) - \mathbf{P}_i^{\text{norm}}(t - \Delta t)}{\Delta t}
\label{eq:joint_velocity_def}
\end{equation}

While mid-air gesture recognition employs filtering techniques to remove jittering \cite{Kumar2026Hybrid}, this cannot be applied to touch detection as smoothing filters destroy the 40~ms deceleration spike ($A_y \ll 0$) upon surface contact and introduce 30 to 90~ms of added latency.

\subsection{Planar Homography and Fiducial Tracking}
\label{subsec:lit_homography}
In order to map camera pixel coordinates into the coordinate system of the printed paper, distortion needs to be removed using a $3 \times 3$ planar homography matrix $H$ \cite{LonguetHiggins1985Machine, Pirchheim2011Homography, Wang2020Relative}:
\begin{equation}
\begin{bmatrix}
x_{\text{XML}} w \\
y_{\text{XML}} w \\
w
\end{bmatrix}
=
H \mathbf{P}_{\text{pixel}}
=
\begin{bmatrix}
h_{11} & h_{12} & h_{13} \\
h_{21} & h_{22} & h_{23} \\
h_{31} & h_{32} & h_{33}
\end{bmatrix}
\begin{bmatrix}
x_{\text{pixel}} \\
y_{\text{pixel}} \\
1
\end{bmatrix}
\label{eq:homography_equation_def}
\end{equation}

The projective scaling factor $w$ allows obtaining Cartesian coordinates on paper:
\begin{equation}
x_{\text{XML}} = \frac{h_{11} x_{\text{pixel}} + h_{12} y_{\text{pixel}} + h_{13}}{h_{31} x_{\text{pixel}} + h_{32} y_{\text{pixel}} + h_{33}}, \quad
y_{\text{XML}} = \frac{h_{21} x_{\text{pixel}} + h_{22} y_{\text{pixel}} + h_{23}}{h_{31} x_{\text{pixel}} + h_{32} y_{\text{pixel}} + h_{33}}
\label{eq:homography_cartesian_def}
\end{equation}

A minimum of four corner correspondences between detected fiducial markers and the printed sheet forms an overdetermined linear system $\mathbf{A}_{8 \times 9} \mathbf{h} = \mathbf{0}$, resolved through Singular Value Decomposition (SVD) \cite{Wang2016AprilTag2, Kallwies2020Determining}.

In experimental results provided by Kalaitzakis et al. \cite{Kalaitzakis2021Fiducial}, AprilTag provides a smaller error of corner localization ($0.017\text{ pixels}$) compared to ArUco and ARTag, even at a camera tilt of $75^\circ$, running in under $1.0\text{ ms}$ on standard commodity CPUs without GPU hardware.

\subsection{Deep Sequential Touch Detection Models}
\label{subsec:lit_deep_touch}
A simple velocity threshold approach ($V_y > \theta$) cannot reliably detect touch due to hand hovering and tendon coupling between fingers \cite{Simao2016Unsupervised, Toshpulatov2024RealTime}. Long Short-Term Memory (LSTM) recurrent neural networks \cite{Hochreiter1997LSTM} retain internal memory over time, thus providing the perfect fit for classifying kinetic impact:
\begin{align}
\mathbf{f}_t &= \sigma(\mathbf{W}_f [\mathbf{h}_{t-1}, \mathbf{x}_t] + \mathbf{b}_f) \label{eq:lstm_forget_def} \\
\mathbf{i}_t &= \sigma(\mathbf{W}_i [\mathbf{h}_{t-1}, \mathbf{x}_t] + \mathbf{b}_i) \label{eq:lstm_input_def} \\
\tilde{\mathbf{C}_t} &= \tanh(\mathbf{W}_c [\mathbf{h}_{t-1}, \mathbf{x}_t] + \mathbf{b}_c) \label{eq:lstm_candidate_def} \\
\mathbf{C}_t &= \mathbf{f}_t \odot \mathbf{C}_{t-1} + \mathbf{i}_t \odot \tilde{\mathbf{C}_t} \label{eq:lstm_cell_def} \\
\mathbf{o}_t &= \sigma(\mathbf{W}_o [\mathbf{h}_{t-1}, \mathbf{x}_t] + \mathbf{b}_o) \label{eq:lstm_output_def} \\
\mathbf{h}_t &= \mathbf{o}_t \odot \tanh(\mathbf{C}_t) \label{eq:lstm_hidden_def}
\end{align}

Recurrent models for sequence learning consider joint movements in sliding temporal windows, examining multi-frame sequences rather than a single static image \cite{Nunez2018Convolutional, Ikram2020Skeleton, Doan2022Efficient}. This separates downward descent, physical surface contact, and finger rebound cleanly.

\subsection{System Architecture Sensing Trade-Offs}
\label{subsec:lit_sensing_tradeoffs}
Table~\ref{tab:lit_sensing_tradeoffs} compares the architectural trade-offs across common optical and physical sensing modalities evaluated in virtual input literature.

\begin{table}[htbp]
\centering
\caption{System architecture design trade-offs across sensing modalities.}
\label{tab:lit_sensing_tradeoffs}
\begin{tabular}{l l l l l}
\toprule
\textbf{Architectural Dimension} & \textbf{Monocular RGB + Planar Fiducials} & \textbf{RGB-D / ToF Depth Sensing} & \textbf{Laser Projection + IR} & \textbf{3D Mid-Air Typing (VR)} \\
\midrule
\textbf{Hardware Sensor Cost} & Low (Commodity Webcam) & High (\$150 to \$400 Depth Camera) & Medium (\$80 to \$150 Projector) & High (\$400 to \$3500 Headset) \\
\textbf{Ambient Light Sensitivity} & Moderate (Robust to indoor lux) & High (Sunlight IR interference) & Extreme (Sunlight washes pattern) & Low (Built-in HMD sensors) \\
\textbf{Surface Alignment} & Dynamic (AprilTag Homography) & Fixed Depth Plane Calibration & Fixed Desktop Projection & None (Mid-air tracking) \\
\textbf{Tactile Feedback} & Passive Surface Tactile Stop & Passive Surface Tactile Stop & Passive Surface Tactile Stop & Complete Absence (Fatigue) \\
\textbf{Layout Customization} & High (Printable XML Layouts) & Moderate (Fixed Spatial GUI) & Low (Hardcoded Laser Pattern) & High (Software VR Grids) \\
\textbf{Edge Compute Latency} & Low ($> 60$ FPS on CPU) & Medium (Point Cloud Compute) & Low (Hardware IR Processing) & High (Heavy HMD Pipeline) \\
\bottomrule
\end{tabular}%
\end{table}

As illustrated in the comparison, monocular RGB sensing with planar fiducials strikes a strong balance between low hardware unit cost, passive desk tactile contact without requiring specialized active sensors, and customizable planar layout tracking relative to active depth sensors or projection hardware.

\section{Reflection and Synthesis of Research Gaps}
\label{sec:lit_reflection}

\subsection{Synthesis of Prior Literature Limitations}
\label{subsec:lit_synthesis_limitations}
However, there are some common problems that can be seen from the discussed systems. Firstly, for monocular systems in physical desktops, previous systems could only track one finger per time, maybe an index fingertip \cite{Zhang2001Visual, Srivastava2012RealTime, Posner2012Single, Thomas2013Camera, Ji2018Local, Khare2019QWERTY}. In the event where users move multiple fingers simultaneously, shadow detection will not work because of the overlapping shadows \cite{Posner2012Single, Yue2014Blind}, and 2D pixel detection will not be able to distinguish the depth of each finger that is coming down.

In order to cope with the lack of depth information, previous 2D systems had to impose pauses of $500\text{ to }1000\text{ ms}$ on each keystroke (dwell time) \cite{Khare2019QWERTY, Chen2024Gaze}, thus making typing very slow and disrupting natural interaction flow. Also, deep neural networks require fast GPUs to run and do not work well on commodity CPUs, achieving fewer than 15~FPS \cite{Enkhbat2020HandKey, Ji2018Local}.

A further concern is the lighting-sensitive nature of the problem. Simple room lights, hand shadows, and variations in skin complexions break down the segmentation and skin color detection process \cite{Thomas2013Camera, Yue2014Blind}. In order to overcome this challenge, several researchers resorted to using expensive three-dimensional sensors, infrared projectors, or even wearable devices like gloves with sensors \cite{Lee2019Virtual, Kudale2016RealTime, Toshpulatov2024RealTime, Pilkov2024Markerless}, yet this approach defeats the purpose of using an inexpensive vision system. Finally, virtually all vision keyboards from earlier work use fixed QWERTY keyboards \cite{Maman2023TypeNet, Enkhbat2020HandKey}. They lack custom-printed keyboards with fiducials, nor do they offer mapping of multiple software profiles to one piece of paper.

\subsection{Formal Research Gap Justification}
\label{subsec:lit_gap_justification}
These drawbacks bring into focus the primary problem that the current work aims at solving; it is the construction of an adaptive paper-based macro keyboard by means of solely a regular monocular RGB camera and a commodity CPU. The key requirements include the detection of multi-touch hits without any additional dwell times, utilization of AprilTag planar homography ($H$) for tracking the paper's position, and providing users with options for defining different software actions on a single printed layout.

This can be done by integrating the techniques of layout design using visualization, planar homography, unit-independent scaling ($L_{\text{hand}}$), direct feature passing, and a light-weighted PyTorch LSTM sequence model, as described in the subsequent chapters.

\section{Summary of the Literature and Research Foundations}
\label{sec:lit_chapter_summary}
This chapter provided an overview of the prior body of knowledge in the areas of virtual input devices, hand pose tracking, planar fiducials, and deep learning touch detection. Past studies on input devices technologies, typing effects on biomechanics, and ergonomics measurement were reviewed. Comparative analysis was conducted by comparing eight sensing approaches using a comparative matrix, whereas the MediaPipe algorithm, scale normalization, AprilTag homography, and PyTorch LSTM sequence classification algorithms were discussed. Finally, the six research questions answered in this thesis were summarized. These principles serve as direct inputs for the methodology and architecture presented in Chapter~\ref{ch:methodology}.
